% ---------------------------------------------------------------------------
% Cluster: computational illustration of the corrected-grid theorem, and three
% results that the experiments motivated:
%   (1) growth certificates for box QPs and where nonconvexity can live;
%   (2) sharpness of the localization radius (filtering alone leaves order
%       sqrt(n kappa) labels per coordinate; grading is what gives log n);
%   (3) localized exact acceptance (exact output without the height bound).
% Labels are prefixed with "exp:". References with prefix core: point to the
% core results (interpolation, filtering, trial invariants, localization).
% ---------------------------------------------------------------------------

\subsection*{Conventions for this part}

Throughout, $F(x)=k+b^{\mathsf T}x+\tfrac12x^{\mathsf T}Hx$ is a quadratic with
symmetric $H$ on $X=\prod_iX_i$, where $X_i=[l_i,u_i]$ for $i\in I_C$ and
$X_i=[l_i,u_i]\cap\Z$ for $i\in I_Z$, with $l_i<u_i$ (fixed coordinates are
substituted). We write $\bar X=\prod_i[l_i,u_i]$ for the continuous hull,
$\ell(x)=\nabla F(x)=Hx+b$, and $d=x-\hat x$ for a reference point $\hat x$.
The coordinate curvature bounds are $L_i=\max\{H_{ii},0\}$ and
$L=\max_iL_i$; along every coordinate line of $\bar X$ the function
$t\mapsto F(\dots,t,\dots)-L_it^2/2$ is concave (it is affine or concave).
In the grid method a coordinate with $L_i>0$ receives the graded grid of the
main text with corrections $d_i(v)=L_iw_i(v)^2/8$; a coordinate with $L_i=0$
receives the two-point grid consisting of the endpoints of its current
domain and correction $0$. Corollary~\ref{exp:cor:eventual} uses this
per-coordinate form; with one common $L$ every coordinate is graded, and
nothing else in this part changes.

A \emph{stage} works in a current box $B^\circ=\prod_iB^\circ_i\subseteq\bar X$
whose integer coordinates have integral endpoints, with grids $G_i\subseteq
B^\circ_i$ containing the endpoints of $B^\circ_i$, corrected objective
$Q=F-D$, minimum $\beta$, min-marginals $m_i(v)$, incumbent value $U$ (the
value of a feasible point), and retained hull $B=\prod_iB_i\subseteq B^\circ$.
We use the following facts from the main text.
\begin{enumerate}
\item[(I)] \emph{Interpolation} (Lemma~\ref{core:lem:interp}). For every
$x\in B^\circ\cap X$ there is a random grid point $Y$ with $\E Q(Y)\le F(x)$
and $Y_i\in\{a,b\}$ whenever $x_i$ lies in the grid interval $[a,b]$ of
$G_i$; if $x_i\in G_i$ then $Y_i=x_i$. With $L_i=0$ and the two-point grid this
holds by concavity along coordinate $i$ (Jensen's inequality), with zero
correction.
\item[(F)] \emph{Filtering history} (Proposition~\ref{core:prop:filter}). If the
stage is reached from the original box by filtering steps with nonincreasing
incumbent values, then every $x\in X$ with $F(x)\le U$ lies in $B$.
\end{enumerate}
Under a unique minimizer $x^*$ with growth constant $g$,
$\kappa=\max\{1,L/g\}$, and a slope with $\theta^2\le1/(8\kappa)$, stage $j$
of a trial (mesh $h_j=s2^{-j}$, center $c=y_{j-1}$, corrected minimizer $y_j$)
satisfies:
\begin{enumerate}
\item[(G1)] $D(z)\le\frac L4nh_j^2+\frac{L\theta^2}2(\norm{z-x^*}^2+\norm{c-x^*}^2)$
for every grid point $z$ (inequality~\eqref{core:eq:energy});
\item[(G2)] $\norm{y_j-x^*}^2\le\kappa nh_j^2$ and $\norm{c-x^*}^2\le4\kappa nh_j^2$
(Lemma~\ref{core:lem:inv}(ii),(iii));
\item[(G3)] $U-\OPT\le F(y_j)-\OPT\le D(y_j)\le\frac78Lnh_j^2$
(Lemma~\ref{core:lem:inv}(iv));
\item[(G4)] for $i\in I_C$, $B_i\subseteq[(y_j)_i-5\sqrt{n\kappa}\,h_j,\,
(y_j)_i+5\sqrt{n\kappa}\,h_j]$; for $i\in I_Z$ the radius is
$1+5\sqrt{n\kappa}\,h_j$ (Lemma~\ref{core:lem:states}).
\end{enumerate}
These are the trial invariants written for one mesh $h_j=s2^{-j}$ in all
coordinates and corrections $d_i\le L(h_j+\theta|z_i-c_i|)^2/8$; their
proofs use only these two facts, and they hold there with the sharper
constants $\frac9{16}$ in (G3) and $4.2$ in (G4). We use the weaker
constants stated. For a two-point coordinate ($L_i=0$) the correction
vanishes, so (G1)--(G3) hold verbatim; (G4) is not used for such
coordinates.

% ---------------------------------------------------------------------------
\subsection{Quadratic growth on a box: a certificate and a necessary condition}

\begin{lemma}[Growth certificates for box QPs]\label{exp:lem:growth}
Let $X=\bar X$ be continuous, $x^*\in X$, $\ell=\ell(x^*)$,
$S=\{i:l_i<x^*_i<u_i\}$ and $A=\{i:x^*_i\in\{l_i,u_i\}\}$.
\begin{enumerate}
\item[(a)] Suppose $\ell_i=0$ for $i\in S$, $\ell_i\ge0$ if $x_i^*=l_i$, and
$\ell_i\le0$ if $x^*_i=u_i$. Let $M=\operatorname{diag}(\mu)$ with
$\mu_i=|\ell_i|/(u_i-l_i)$ for $i\in A$ and $\mu_i=0$ for $i\in S$. Then
\begin{equation}\label{exp:eq:minorant}
F(x)-F(x^*)\ge\tfrac12(x-x^*)^{\mathsf T}(H+2M)(x-x^*)\qquad(x\in X).
\end{equation}
Consequently, if $H+2M\succeq2\gamma I$ for some $\gamma>0$, then $x^*$ is the
unique minimizer and $F(x)-\OPT\ge\gamma\norm{x-x^*}^2$ on $X$.
\item[(b)] Conversely, if $x^*$ is a minimizer and
$F(x)-\OPT\ge g\norm{x-x^*}^2$ on $X$ with $g>0$, then $\ell_S=0$ and
$H_{SS}\succeq2gI$. In particular $g\le\frac12v^{\mathsf T}Hv/\norm v^2$ for
every nonzero $v$ supported on $S$; if $S\ne\emptyset$ then every valid
curvature bound satisfies $L\ge2g$, so $\kappa\ge2$; and if $S=\{1,\dots,n\}$
then $H\succeq2gI$, so $F$ is strongly convex.
\end{enumerate}
\end{lemma}

\begin{proof}
(a) Let $d=x-x^*$. Since $F$ is quadratic,
$F(x)-F(x^*)=\ell^{\mathsf T}d+\frac12d^{\mathsf T}Hd$ exactly. For $i\in S$,
$\ell_id_i=0$. For $i\in A$ with $x_i^*=l_i$ we have $0\le d_i\le u_i-l_i$
and $\ell_i\ge0$, so $\ell_id_i=|\ell_i|d_i\ge|\ell_i|d_i^2/(u_i-l_i)=\mu_id_i^2$;
the case $x_i^*=u_i$ is symmetric. Summing gives~\eqref{exp:eq:minorant}. If
$H+2M\succeq2\gamma I$, the right side is at least $\gamma\norm d^2$, which is
positive for $d\ne0$.

(b) For $i\in S$ both directions $\pm e_i$ are feasible from $x^*$, so
minimality gives $\ell_i=0$. Let $v\ne0$ be supported on $S$. For all
sufficiently small $|t|$, $x^*+tv\in X$, and
$F(x^*+tv)-\OPT=\frac{t^2}2v^{\mathsf T}Hv\ge gt^2\norm v^2$. Hence
$H_{SS}\succeq2gI$. Taking $v=e_i$ with $i\in S$ gives $H_{ii}\ge2g$; every
valid curvature bound satisfies $L\ge H_{ii}$, because
$t\mapsto F(\dots,t,\dots)-Lt^2/2$ has second derivative $H_{ii}-L$. If
$S$ is everything, $H=H_{SS}\succeq2gI$.
\end{proof}

\begin{remark}[Where nonconvexity can live]\label{exp:rem:nonconvex}
Part~(b) says that for a continuous box QP, quadratic growth at an interior
minimizer forces strong convexity, and in general the Hessian block of the
free coordinates must be positive definite. A nonconvex continuous instance
covered by the growth theorem therefore has its negative curvature in
directions that involve coordinates at bounds, where first-order terms
$\ell_id_i$ provide the growth, or in integer coordinates. Part~(a) is a
sufficient certificate of exactly this type. For an integer coordinate one
may also use $\ell_id_i\ge-|\ell_i|d_i^2$, valid because $|d_i|\le d_i^2$
for $d_i\in\Z$; this replaces $\mu_i$ by $-|\ell_i|$. The planted instances of
the computational illustration are built from part~(a): $H_{SS}$ has a
prescribed smallest eigenvalue, the entries touching $A$ make $H$ indefinite,
and $\mu$ is chosen so that $H+2M\succeq2\gamma I$ is verified exactly; then
$\gamma\le g\le\frac12v^{\mathsf T}H_{SS}v/\norm v^2$ for any rational $v$,
which brackets $\kappa$. Such instances are benign in one sense: once $x^*$
is known, \eqref{exp:eq:minorant} certifies it directly.
\end{remark}

% ---------------------------------------------------------------------------
\subsection{Filtered uniform grids}

With uniform grids (slope $0$) the analysis needs no slope condition and no
knowledge of $g$: every correction is at most $Lh^2/8$, wherever the grid
point lies.

\begin{proposition}[Filtered uniform grids]\label{exp:prop:uniform}
Assume a unique minimizer $x^*$ with growth constant $g$. Consider any stage
whose grid intervals all have length at most $h$ in every coordinate with
$L_i>0$ (as for the uniform rule with mesh $h$), with corrected minimizer $y$
and incumbent value $U\le F(y)$. Then
$\norm{y-x^*}\le\sqrt{n\kappa/8}\,h$, and for every continuous coordinate
with $L_i>0$ the retained hull lies within distance
\[
\Bigl(\bigl(\tfrac12+\tfrac1{2\sqrt2}\bigr)\sqrt{n\kappa}+1\Bigr)h
\]
of $y_i$. For the uniform rule the next stage therefore has at most
$(2+\sqrt2)\sqrt{n\kappa}+7$ labels in each such coordinate, at every
stage and for every accuracy, without any condition on $\kappa$.
\end{proposition}

\begin{proof}
Every node $v$ has $d_i(v)\le L_ih^2/8\le Lh^2/8$, so $D(z)\le nLh^2/8$
for every grid point $z$. By (I), $\beta\le\OPT$, so
$F(y)-\OPT\le F(y)-\beta=D(y)\le nLh^2/8$, and growth gives
$\norm{y-x^*}^2\le nLh^2/(8g)\le n\kappa h^2/8$. A retained interval has an
endpoint $v$ with $m_i(v)\le U$; let $z$ be a grid point with $z_i=v$ and
$Q(z)=m_i(v)$. Then
\[
g\norm{z-x^*}^2\le F(z)-\OPT=Q(z)+D(z)-\OPT\le F(y)-\OPT+D(z)
\le\frac{nLh^2}4,
\]
so $|v-x^*_i|\le\sqrt{n\kappa/4}\,h$ and $|v-y_i|\le
(\frac12+\frac1{2\sqrt2})\sqrt{n\kappa}\,h$. The interval has length at most
$h$, which gives the radius, and the hull of such intervals obeys the same
bound. The next uniform grid is centered at $y$ with spacing $h/2$; its
nodes are the points of $y_i+\frac h2\Z$ in the hull, at most
$4\cdot\frac Rh+1$ of them for radius $R$, plus the two hull endpoints. This
gives at most $4((\frac12+\frac1{2\sqrt2})\sqrt{n\kappa}+1)+3
=(2+\sqrt2)\sqrt{n\kappa}+7$ labels.
\end{proof}

% ---------------------------------------------------------------------------
\subsection{The localization radius is attained}

Lemma~\ref{core:lem:states} bounds the retained radius by a constant times
$\sqrt{n\kappa}\,h_j$ ($5\sqrt{n\kappa}\,h_j$ in (G4)). The following
instance shows that order $\sqrt{n\kappa}\,h_j$ is attained
for every $\kappa\ge2$, with exact coordinate curvature. The reason is that
the slack of the lower bound, $U-\beta$, is a sum of $n$ correction terms,
whereas a single coordinate is removed only when its own excess exceeds the
whole slack.

\begin{proposition}[Sharpness of localization]\label{exp:prop:sharp}
Let $m\ge1$, $n=2m$, $0\le\sigma<\Lambda$, and on $\bar X=[-1,1]^n$ with
coordinates $(u_1,v_1,\dots,u_m,v_m)$ let
\[
F(x)=\sum_{k=1}^m\Bigl[\frac\Lambda2(u_k^2+v_k^2)+\sigma u_kv_k\Bigr].
\]
Then $x^*=0$, $\OPT=0$, every coordinate has exact curvature $L=\Lambda$, the
best growth constant is $g=(\Lambda-\sigma)/2$, and
$\kappa=L/g=2\Lambda/(\Lambda-\sigma)$ takes every value in $[2,\infty)$.
Consider a stage whose box contains $[-\rho,\rho]^n$ and whose grids
$G_i$ all contain $0$, with both grid intervals at $0$ of length at least $h$.
Then for every incumbent value $U\ge0$, every coordinate $i$ and every node
$a\in G_i$ with
\[
|a|\le\min\Bigl\{\rho,\ h\sqrt{(n-2)\kappa/16}\Bigr\}
\]
we have $m_i(a)\le U$. Hence every grid interval with such an endpoint is
retained.
\end{proposition}

\begin{proof}
The Hessian is block diagonal with blocks
$\bigl(\begin{smallmatrix}\Lambda&\sigma\\\sigma&\Lambda\end{smallmatrix}\bigr)$,
whose eigenvalues are $\Lambda\pm\sigma>0$. Hence $F(x)\ge\frac{\Lambda-\sigma}2
\norm x^2$ with equality along $u_k=-v_k$, so $x^*=0$, $\OPT=0$ and
$g=(\Lambda-\sigma)/2$. Each diagonal entry is $\Lambda$. Since
$\kappa=2\Lambda/(\Lambda-\sigma)$ increases continuously from $2$ to
$\infty$ as $\sigma$ goes from $0$ to $\Lambda$, every $\kappa\ge2$ occurs.

The corrected objective separates over pairs,
$Q=\sum_kq_k(u_k,v_k)$ with
$q_k(u,v)=\frac\Lambda2(u^2+v^2)+\sigma uv-\frac\Lambda8(w(u)^2+w(v)^2)$,
where $w$ denotes the largest adjacent interval length in the relevant
grid. Therefore the min-marginal of $u_1$ is
$m_{u_1}(a)=\min_vq_1(a,v)+\sum_{k\ge2}\min q_k$, the minima taken over the
grids. By hypothesis $w(0)\ge h$ in every coordinate, so
$\min q_k\le q_k(0,0)\le-\Lambda h^2/4$ for each $k\ge2$, and
$\sum_{k\ge2}\min q_k\le-(m-1)\Lambda h^2/4=-(n-2)\Lambda h^2/8$.

Fix a node $a$ with $|a|\le\rho$ and put $t^*=-\sigma a/\Lambda$; then
$|t^*|\le|a|\le\rho$, so $t^*$ lies in the box of $v_1$ and in some grid
interval of $G_{v_1}$ of length $\lambda$, say. Let $v_a$ be the endpoint of
that interval nearest to $t^*$. Then $|v_a-t^*|\le\lambda/2$ and
$w(v_a)\ge\lambda$. Completing the square,
$\frac\Lambda2v^2+\sigma av=\frac\Lambda2(v-t^*)^2-\frac{\sigma^2a^2}{2\Lambda}$,
so
\[
q_1(a,v_a)\le\frac{\Lambda^2-\sigma^2}{2\Lambda}a^2
+\frac\Lambda2\Bigl(\frac\lambda2\Bigr)^2-\frac\Lambda8\lambda^2
=\frac{\Lambda^2-\sigma^2}{2\Lambda}a^2 ,
\]
where we also dropped the nonpositive term $-\frac\Lambda8w(a)^2$. Since
$\Lambda-\sigma=2g$ and $\Lambda+\sigma=2(\Lambda-g)$,
$(\Lambda^2-\sigma^2)/(2\Lambda)=2g(\Lambda-g)/\Lambda$. Thus
\[
m_{u_1}(a)\le\frac{2g(\Lambda-g)}{\Lambda}a^2-\frac{(n-2)\Lambda h^2}8,
\]
which is at most $0\le U$ when
$a^2\le(n-2)\Lambda^2h^2/(16g(\Lambda-g))=(n-2)h^2\kappa^2/(16(\kappa-1))$.
Because $\kappa^2/(\kappa-1)\ge\kappa$, the hypothesis
$a^2\le(n-2)\kappa h^2/16$ suffices. The coordinates $v_k$ and the other
pairs are symmetric. An interval with an endpoint $a$ satisfying
$m_i(a)\le U$ passes the filtering test.
\end{proof}

\begin{corollary}[Filtered uniform grids need order $\sqrt{n\kappa}$ labels]
\label{exp:cor:uniform}
In the setting of Proposition~\ref{exp:prop:sharp} with $n\ge4$, run the
filtered algorithm with uniform grids (slope $\theta=0$), started at the
minimizer $0$. Let $r=\lfloor\sqrt{(n-2)\kappa/16}\rfloor$. At every stage the
corrected minimizer is $0$; at every stage $j\ge1$ with $(r+1)h_j\le1$ the
retained box contains $[-(r+1)h_j,(r+1)h_j]^n$, and the grid of stage $j+1$
has at least $4r+5\ge\sqrt{(n-2)\kappa}+1$ nodes in every coordinate. For
graded grids with slope $\theta>0$: if the box of a stage $j$ contains
$[-R,R]^n$, where $R=h_j\sqrt{(n-2)\kappa/16}$, and its grids contain $0$
with $w_i(0)\ge h_j$, then every coordinate of the grid of stage $j+1$ has
at least $1+\log(1+2\theta R/h_j)/\log(1+\theta)$ nodes.
\end{corollary}

\begin{proof}
Uniform rule, center. Suppose a stage is centered at $0$ with box
$B^\circ\ni0$. Its grid in coordinate $i$ consists of the multiples of $h_j$
in $B^\circ_i$ and the endpoints of $B^\circ_i$. All grid intervals have length
at most $h_j$, and the longest one adjacent to $0$ is at least as long as any
other interval, because only the outermost intervals are shortened by the
endpoints. Hence $w_i(v)\le w_i(0)$ for all nodes, so for every grid point
$x\ne0$,
$Q(x)-Q(0)\ge F(x)-F(0)\ge\frac{\Lambda-\sigma}2\norm x^2>0$. The corrected
minimizer is $0$, and so is the next center; by induction every stage is
centered at $0$.

Uniform rule, retained box. Put $\rho_j=\min\{1,2(r+1)h_j\}$; we show by
induction that the box $B_j$ of stage $j$ contains $[-\rho_j,\rho_j]^n$.
This holds for $j=0$ ($B_0=[-1,1]^n$). At stage $0$ ($h_0=2$) the grid is
$\{-1,0,1\}$ in every coordinate, with $w_i(0)=1$; by
Proposition~\ref{exp:prop:sharp} with $h=1$ the node $0$ passes, so both
intervals are retained and $B_1=[-1,1]^n\supseteq[-\rho_1,\rho_1]^n$. Let
$j\ge1$. Then $h_j\le1$, hence $h_j\le\rho_j$, so $\pm h_j$ are nodes and
$w_i(0)=h_j$. Proposition~\ref{exp:prop:sharp} with $h=h_j$ and
$\rho=\rho_j$ retains every interval with an endpoint $kh_j$, where
$|k|\le r$ and $|k|h_j\le\rho_j$. If $(r+1)h_j\le1$, then $(r+1)h_j\le\rho_j$
and the retained intervals $[kh_j,(k+1)h_j]$, $-r-1\le k\le r$, give
$B_{j+1}\supseteq[-(r+1)h_j,(r+1)h_j]^n=[-\rho_{j+1},\rho_{j+1}]^n$. If
$(r+1)h_j>1$, then $\rho_j=1$, $B_j=[-1,1]^n$, and the retained intervals
reach the box boundary (the outermost ones possibly shortened), so
$B_{j+1}=[-1,1]^n\supseteq[-\rho_{j+1},\rho_{j+1}]^n$. This completes the
induction. When $(r+1)h_j\le1$, the grid of stage $j+1$ has spacing $h_j/2$
and contains every multiple of $h_j/2$ in $[-(r+1)h_j,(r+1)h_j]$, that is,
$4(r+1)+1$ nodes; and $r+1\ge\sqrt{(n-2)\kappa/16}$ gives
$4r+5\ge\sqrt{(n-2)\kappa}+1$.

Graded rule. Apply Proposition~\ref{exp:prop:sharp} with $h=h_j$ and
$\rho=R$. Every point $t\in[-R,R]$ lies in a grid interval $[a,b]$; if
$t\ge0$ then $0\le a\le R$ (the node $0$ lies in the grid), and if $t<0$ then
$-R\le b\le0$, so the interval has an endpoint of absolute value at most
$R$ and is retained. Hence the retained box contains $[-R,R]$ in every
coordinate. The next grid starts from its center $c'$ inside this box and
uses unclipped steps $h'+\theta t$ with $h'=h_j/2$, so after $k$ steps the
distance from $c'$ is at most $h'((1+\theta)^k-1)/\theta$. One side of $c'$
has length at least $R$ within the box, and must be covered; hence
$k\ge\log(1+\theta R/h')/\log(1+\theta)$ nodes lie on that side, besides
$c'$.
\end{proof}

\begin{remark}[What filtering and grading each contribute]\label{exp:rem:roles}
Proposition~\ref{exp:prop:uniform} and Corollary~\ref{exp:cor:uniform}
show that filtered uniform grids have $\Theta(\sqrt{n\kappa})$ labels per
coordinate at every late stage, between $\sqrt{(n-2)\kappa}+1$ (for the pair
instance) and $(2+\sqrt2)\sqrt{n\kappa}+7$ (always), independent of the
accuracy, with no slope condition and no growth constant in the algorithm:
filtering alone removes the accuracy dependence, and a single run suffices. With bags of size
$p$ the tables then have order $(n\kappa)^{p/2}$ states, so the exponent of
$n$ grows with $p$; this is polynomial for fixed $p$ but not of the form
$f(p,\kappa)\poly(n)$. Graded grids replace $\sqrt{n\kappa}$ by
$\theta^{-1}\log(1+\theta\sqrt{n\kappa})$, and with $\theta$ of order
$\kappa^{-1/2}$ by order $\sqrt\kappa\log n$, which is what makes
$\lceil\log_2(n+2)\rceil^p\le(C_0p)^p(n+2)$ applicable. The same proposition
shows that the localization radius of Lemma~\ref{core:lem:states} is sharp up
to an absolute factor (at most $20\sqrt2$ for $n\ge4$ with the radius
$5\sqrt{n\kappa}\,h_j$ of (G4)). The phenomenon is the
grid analogue of the cluster problem of branch-and-bound
\cite{DuKearfott1994,WechsungSchaberBarton2014,KannanBarton2017}: a
second-order lower bound whose error constant grows with the number of
coordinates leaves a neighborhood of the minimizer whose size grows with
$n$.
\end{remark}

% ---------------------------------------------------------------------------
\subsection{Exact acceptance from a localized certificate}

The exact-output theorem accepts a rational candidate once the certified gap
is below $1/(VW)$, where $V$ bounds the height of the optimal value. On
small instances this already requires gaps near $2^{-100}$ to $2^{-350}$
(Section~\ref{exp:sec:computation}). The next result accepts a candidate as
soon as filtering has localized it, using only first- and second-order
information on the retained box.

\begin{lemma}[Node exclusion]\label{exp:lem:node}
In a stage with box $B^\circ$, let $i$ be a coordinate and $v\in G_i$. Then
$F(x)\ge m_i(v)$ for every $x\in B^\circ\cap X$ with $x_i=v$. Consequently,
for every $x\in B^\circ\cap X$ with $F(x)\le U$:
\begin{enumerate}
\item[(a)] if $i\in I_Z$, then $x_i$ belongs to the set $K_i$ consisting of the
nodes $v\in G_i$ with $m_i(v)\le U$ and the integers strictly inside grid
intervals $[a,b]$ with $b-a\ge2$ and $\min\{m_i(a),m_i(b)\}\le U$;
\item[(b)] if $L_i=0$, $G_i=\{a,b\}$ consists of the endpoints of $B^\circ_i$,
and $m_i(a)>U$, then the point $x'$ obtained from $x$ by setting $x'_i=b$
satisfies $x'\in B^\circ\cap X$ and $F(x')\le F(x)$.
\end{enumerate}
\end{lemma}

\begin{proof}
If $x_i=v\in G_i$, the random grid point of (I) has $Y_i=v$ almost surely,
so $m_i(v)\le Q(Y)$ pointwise and $m_i(v)\le\E Q(Y)\le F(x)$.
(a) If $x_i\in G_i$ this gives $m_i(x_i)\le U$. Otherwise $x_i$ is an
integer strictly inside a grid interval $[a,b]$, so $b-a\ge2$, and (I) gives
$\min\{m_i(a),m_i(b)\}\le F(x)\le U$.
(b) Since $L_i=0$, $t\mapsto F(x_1,\dots,t,\dots,x_n)$ is concave on
$[a,b]$, so its minimum over $[a,b]$ is attained at $a$ or $b$, and
$F(x)\ge\min\{F(x|_{x_i=a}),F(x|_{x_i=b})\}$. The first value is at least
$m_i(a)>U\ge F(x)$, so the minimum is $F(x')$ with $x'_i=b$. The endpoint
$b$ is feasible, and integral when $i\in I_Z$.
\end{proof}

\begin{proposition}[Localized exact acceptance]\label{exp:prop:local}
Consider a stage with incumbent value $U$ and retained hull $B$ satisfying
(F). Define a box $B'\subseteq B$ coordinatewise: for $i\in I_Z$ with
$L_i>0$ let $B'_i$ be the hull of $K_i\cap B_i$ (Lemma~\ref{exp:lem:node}(a));
for $i$ with $L_i=0$ and two-point grid $\{a,b\}$ let $B'_i=\{b\}$ if
$m_i(a)>U\ge m_i(b)$, $B'_i=\{a\}$ if $m_i(b)>U\ge m_i(a)$, and $B'_i=B_i$
otherwise; let $B'_i=B_i$ for the remaining coordinates. Let
$\hat x\in B'\cap X$ satisfy $F(\hat x)\le U$, let $\ell=\ell(\hat x)$,
$W=\{i:B'_i\text{ has positive length}\}$, and $B'_i=[\alpha_i,\gamma_i]$.
Suppose every $i\in W$ falls under one of the cases
\[
\begin{array}{lll}
\text{(i)}&\hat x_i=\alpha_i,\ \ell_i\ge0:&\mu_i=\ell_i/(\gamma_i-\alpha_i);\\
\text{(ii)}&\hat x_i=\gamma_i,\ \ell_i\le0:&\mu_i=-\ell_i/(\gamma_i-\alpha_i);\\
\text{(iii)}&\alpha_i<\hat x_i<\gamma_i,\ \ell_i=0:&\mu_i=0;\\
\text{(iv)}&i\in I_Z\ \text{otherwise}:&\mu_i=-|\ell_i|,
\end{array}
\]
and that $H_{WW}+2\operatorname{diag}(\mu_W)\succeq0$. Then $F(\hat x)=\OPT$.
\end{proposition}

\begin{proof}
Let $x\in X$. If $F(x)>U$ then $F(x)>U\ge F(\hat x)$. Otherwise $x\in B$ by
(F). Apply Lemma~\ref{exp:lem:node}(b) successively to the coordinates with
$L_i=0$ whose $B'_i$ is a single endpoint; each step keeps the point in
$B\cap X$ (the replaced coordinate stays in its interval) and does not
increase $F$, so each intermediate point still has value at most $U$, and the
lemma applies again. The result $x''\in B\cap X$ has $F(x'')\le F(x)\le U$
and agrees with $B'_i$ on those coordinates. By Lemma~\ref{exp:lem:node}(a)
applied to $x''$, every integer coordinate with $L_i>0$ satisfies
$x''_i\in K_i\cap B_i\subseteq B'_i$. Hence $x''\in B'\cap X$.

Let $d=x''-\hat x$. Both points lie in $B'$, so $d_i=0$ for $i\notin W$, and
for $i\in W$ we have $\alpha_i-\hat x_i\le d_i\le\gamma_i-\hat x_i$. We
claim $\ell_id_i\ge\mu_id_i^2$ for $i\in W$. In case (i),
$0\le d_i\le\gamma_i-\alpha_i$ and $\ell_id_i=\ell_i|d_i|\ge
\ell_id_i^2/(\gamma_i-\alpha_i)$; case (ii) is symmetric; in case (iii) both
sides vanish; in case (iv) $d_i\in\Z$, so $|d_i|\le d_i^2$ and
$\ell_id_i\ge-|\ell_i||d_i|\ge-|\ell_i|d_i^2$. Therefore
\[
F(x'')-F(\hat x)=\ell^{\mathsf T}d+\tfrac12d^{\mathsf T}Hd\ge
\tfrac12d_W^{\mathsf T}\bigl(H_{WW}+2\operatorname{diag}(\mu_W)\bigr)d_W\ge0 .
\]
Thus $F(x)\ge F(x'')\ge F(\hat x)$ for every $x\in X$, and $\hat x$ is
feasible.
\end{proof}

The test needs one exact $LDL^{\mathsf T}$ factorization of a $|W|\times|W|$
matrix, which an independent checker can repeat after replaying the
filtering history. The candidate is arbitrary; in practice it is the
stationary point of the face on which a feasible incumbent lies. Validity
needs neither uniqueness nor growth. The next corollary shows that the test
eventually succeeds under growth and strict complementarity, after a number of
stages governed by a geometric margin instead of the height bound.

\begin{corollary}[Acceptance after finitely many stages]\label{exp:cor:eventual}
Let $x^*$ be the unique minimizer with growth constant $g$, and use the
per-coordinate grids of the conventions. Put $C=\{i:L_i=0\}$,
$Z=I_Z\setminus C$, $S=\{i\in I_C:l_i<x^*_i<u_i\}$, and
$A=\{i\in I_C\setminus C:x^*_i\in\{l_i,u_i\}\}$. Assume strict
complementarity: $\ell_i(x^*)\ne0$ for all $i\in A$. Let
$\lambda_A=\min_{i\in A}|\ell_i(x^*)|$,
$\Gamma=\norm{H_{AA}}_2+\norm{H_{SA}}_2^2/(2g)$,
$\delta_C=\min_{i\in C}(u_i-l_i)$, and
\[
h^*=\frac1{\sqrt{n\kappa}}\min\Bigl\{\tfrac12\ [Z\ne\emptyset],\ \
\tfrac45\,\delta_C\ [C\ne\emptyset],\ \ \frac{\lambda_A}{5\Gamma}\
[A\ne\emptyset]\Bigr\},
\]
where a bracketed term is omitted when its condition fails (and
$h^*=\infty$ if all are omitted). In a trial with $\theta^2\le1/(8\kappa)$,
at every stage $j$ with $h_j\le h^*$ the test of
Proposition~\ref{exp:prop:local} accepts $\hat x=x^*$. In particular it
accepts after at most $\lceil\log_2(s/h^*)\rceil+1$ stages of that trial.
\end{corollary}

\begin{proof}
Note that $S\cap C=\emptyset$: for $i\in S$, Lemma~\ref{exp:lem:growth}(b)
gives $H_{ii}\ge2g>0$. Similarly, for $i\in C$ the function
$t\mapsto F(x^*+te_i)$ is concave along the coordinate line through $x^*$
and $F(x^*+te_i)-\OPT\ge gt^2>0$ for feasible $t\ne0$, so $x^*_i$ is an
endpoint of $X_i$; also $B_i=B^\circ_i=X_i$, because a coordinate with a
single grid interval is never narrowed. Fix a stage $j$ with $h_j\le h^*$.

\emph{Node exclusion.} Let $i\in Z\cup C$ and let $v\ne x^*_i$ be a node of
$G_i$; put $\delta=|v-x^*_i|$, so $\delta\ge1$ for $i\in Z$ and
$\delta=u_i-l_i\ge\delta_C$ for $i\in C$. Let $z$ be a grid point with
$z_i=v$ and $Q(z)=m_i(v)$. By growth and (G1), (G2), and
$L\theta^2/2\le g/16$,
\[
m_i(v)-\OPT\ge g\norm{z-x^*}^2-D(z)\ge\tfrac{15}{16}g\delta^2
-\Bigl(\frac L4+\frac{g\kappa}4\Bigr)nh_j^2 .
\]
With (G3) and $L\le\kappa g$ this gives
$m_i(v)-U\ge g\bigl(\frac{15}{16}\delta^2-\frac{11}8\kappa nh_j^2\bigr)>0$
whenever $h_j^2<\frac{15}{22}\delta^2/(\kappa n)$, which holds since
$h_j\le h^*$ and $\frac12<\frac45<\sqrt{15/22}$. Hence for $i\in C$
the rule sets $B'_i=\{x^*_i\}$ (the node $x^*_i$ is never excluded, because
$m_i(x^*_i)\le\OPT\le U$ by Lemma~\ref{exp:lem:node}). For $i\in Z$, all
nodes other than $x^*_i$ have $m_i(v)>U$. Moreover, by (G2),
$|c_i-x^*_i|\le2\sqrt{\kappa n}\,h_j\le1$, and every
integer step taken within distance $5$ of $c_i$ equals
$\max\{1,\lfloor h_j+\theta t\rfloor\}=1$ because
$h_j+\theta t\le\frac12+\frac54<2$. So the grid contains every integer of
$B^\circ_i$ within distance $4$ of $x^*_i$, all grid intervals meeting
$[x^*_i-4,x^*_i+4]$ have length one, and by (G4) and (G2) the hull $B_i$
lies within $1+6\sqrt{n\kappa}\,h_j\le4$ of $x^*_i$. Thus $K_i\cap B_i=\{x^*_i\}$
and $B'_i=\{x^*_i\}$.

\emph{The remaining coordinates.} Hence $W\subseteq S\cup A$, and
$x^*\in B'$ because $x^*\in B$ by (F). For $i\in S\cap W$ we have
$\ell_i(x^*)=0$, so case (i), (ii) or (iii) applies with $\mu_i=0$. For
$i\in A\cap W$, $x^*_i$ is an endpoint of $X_i$, hence of $B'_i=B_i$, and
the first-order conditions give the sign required by case (i) or (ii); by
(G4) the length of $B'_i$ is at most $10\sqrt{n\kappa}\,h_j$, so
$\mu_i\ge\lambda_A/(10\sqrt{n\kappa}\,h_j)\ge\Gamma/2$.

\emph{Positive semidefiniteness.} Let $A'=A\cap W$ and
$N=H_{WW}+2\operatorname{diag}(\mu_W)$. If $A'=\emptyset$, then
$N=H_{SS}\succeq2gI$ by Lemma~\ref{exp:lem:growth}(b) (all of $S$ lies in
$W$, since a continuous coordinate with $L_i>0$ keeps a hull of positive
length). Otherwise $H_{SS}\succ0$, and $N\succeq0$ if the Schur complement
$H_{A'A'}+2\operatorname{diag}(\mu_{A'})-H_{A'S}H_{SS}^{-1}H_{SA'}$ is
positive semidefinite. Its smallest eigenvalue is at least
$2\min_{A'}\mu_i-\norm{H_{AA}}_2-\norm{H_{SA}}_2^2/(2g)\ge0$, using
$\norm{H_{SS}^{-1}}_2\le1/(2g)$ and that submatrices have smaller norms.
Finally $F(x^*)=\OPT\le U$, so all hypotheses of
Proposition~\ref{exp:prop:local} hold. The stage count follows from
$h_j=s2^{-j}$.
\end{proof}

\begin{remark}[Comparison with the height rule]\label{exp:rem:height}
For rational data the margins are not exponentially worse than the height
bound. By the rational-height lemma the unique minimizer $x^*$ has a common
coordinate denominator at most $R$. If $D'$ is a common denominator of the
entries of $H$, $b$ and the bounds, each nonzero $\ell_i(x^*)=(Hx^*+b)_i$ has
denominator dividing $D'R$, so $\lambda_A\ge1/(D'R)$; also
$\delta_C\ge1/D'$, and $\log\Gamma=O(I+\log\kappa)$ because $g\ge L/\kappa$. Hence $\log_2(s/h^*)=\poly(I)+O(\log\kappa)$, as for the
height rule. The difference is practical: the height rule needs a gap below
$1/(VW)$ with $V=DR^2$ computed from a Hadamard-type bound, whereas
Corollary~\ref{exp:cor:eventual} needs only that the retained box be small
compared with the strict-complementarity margin. The idea of certifying a
neighborhood of a candidate by local information is the exclusion-box
technique of interval branch-and-bound \cite{Neumaier2004}; here it is
carried out in exact rational arithmetic on the box produced by the
filtering history. Proposition~\ref{exp:prop:local} remains valid for
nonunique minima, but the test can then fail: with two optimal vertices in a
concave coordinate, neither endpoint is ever excluded.
\end{remark}

\begin{remark}[Node-level filtering of integer labels]\label{exp:rem:integerfilter}
Lemma~\ref{exp:lem:node}(a) also justifies a stronger filtering rule for
integer coordinates: discard every label outside $K_i$. With interval
filtering alone a unit interval $[x^*_i,x^*_i+1]$ is always retained through
its good endpoint, which is the source of the additive $1$ in the integer
case of Lemma~\ref{core:lem:states}. With the node rule a retained label $v$
has $m_i(v)\le U$, and the argument of that lemma bounds $|v-y_i|$ exactly as
for a continuous endpoint; so the integer radius holds without the additive
term.
\end{remark}

% ---------------------------------------------------------------------------
\subsection{Computational illustration (draft for the paper)}\label{exp:sec:computation}

The aim is to illustrate the mechanism of the main theorem, not to rank
solvers. All runs use the exact-rational reference implementation
unchanged; every certificate produced was replayed by the independent
checker and accepted. Timings are single measurements with one numerical
thread on a shared 36-thread Intel Xeon w5-2565X with load average 11--16;
they are indicative only. The full study (290 runs) took 248\,s of wall time
with four worker processes.

\paragraph{Instances.} Planted instances (Remark~\ref{exp:rem:nonconvex})
live on $[-1,1]^n$ with Hessian $2I+C$ supported on a path, a random tree,
or a band of bandwidth $2$ or $3$. About $n/4$ coordinates are active at
$x^*$; $H_{SS}$ has smallest eigenvalue $4/\kappa_{\rm target}$, the other
entries make $H$ indefinite ($\lambda_{\min}(H)$ between $-1.8$ and $-7.2$),
and an exact certificate brackets $\kappa$ within $11\%$; $\OPT=0$. Random
instances (for exact output) have signed sparse Hessians, rational bounds
and one or two integer coordinates, without planted optima.

\paragraph{Accuracy (E1).} Table~\ref{exp:tab:e1} shows states per stage
for four variants with the theorem's slope $\theta=1/8$
($\kappa\in[4.0,4.45]$, medians over three seeds). Both filtered variants
plateau; without filtering the states grow polynomially (graded) or
exponentially (uniform) in the stage index until the cap of $10^5$
states. On the filtered graded runs the measured quantities stayed far
inside the analysis: $\norm{y_j-x^*}^2\le0.047\,\kappa nh_j^2$,
$D(y_j)\le0.155\,Lnh_j^2$ (bound $\frac9{16}$), retained radius at most
$0.135\cdot4.2\sqrt{n\kappa}\,h_j$, and at most $0.0034$ times the label
cap $100\theta^{-1}\lceil\log_2(n+2)\rceil$ used by the implementation. The planted optimizer was never
filtered and every enclosure contained $\OPT$.

\begin{table}[ht]\centering\small
\begin{tabular}{lrrrrrr}
\toprule
& \multicolumn{3}{c}{path, $n=16$, $p=2$} & \multicolumn{3}{c}{band, $n=8$, $p=4$}\\
\cmidrule(lr){2-4}\cmidrule(lr){5-7}
Variant & $j=3$ & $j=7$ & $j=15$ & $j=2$ & $j=3$ & $j=15$\\
\midrule
graded, filtered & 724 & 1{,}014 & 1{,}010 & 3{,}534 & 10{,}895 & 16{,}696\\
uniform, filtered & 916 & 1{,}118 & 1{,}086 & 3{,}966 & 13{,}320 & 16{,}696\\
graded, not filtered & 1{,}047 & 18{,}748 & cap & 4{,}962 & 22{,}857 & cap\\
uniform, not filtered & 1{,}431 & cap & cap & 5{,}658 & 46{,}110 & cap\\
\bottomrule
\end{tabular}
\caption{Bag-table states in stage $j$ (mesh $2^{1-j}$), medians of three
planted instances. ``cap'': the next stage would exceed $10^5$ states.}
\label{exp:tab:e1}
\end{table}

\paragraph{Conditioning (E2) and dimension (E3).} On paths with $n=16$ and
$\kappa_{\rm target}=2,\dots,256$, the plateau label count per coordinate
grows from $13$ to $55$ (uniform), $9$ to $49$ (graded, theorem's $\theta$
from $1/8$ to $1/64$) and $9$ to $33.5$ (graded, $\theta=1/4$); the retained
radius grows like $\sqrt\kappa$ and stays below $0.09\cdot5\sqrt{n\kappa}\,h_j$
with the theorem's $\theta$. With $\theta=1/4$, which violates
$\theta^2\le1/(8\kappa)$ on all these instances, contraction visibly fails
for $\kappa\ge64$: the last-stage gap divided by $Lnh_j^2$ is $4.7$, $15$ and
$61$ at $\kappa=64,128,256$, against about $0.11$ with the theorem's $\theta$.
In the first trial of the default schedule ($\theta=1/4$, $15$ stages) the
retained radius even exceeded the localization bound $5\sqrt{n\kappa}\,h_j$,
by a factor up to $1.45$, so the slope condition is not an artifact of the
proof. The default capped schedule restarted exactly in these cases: its first trial
($\theta=1/4$) succeeded for $\kappa\le32$, the second ($\theta=1/8$) for
$\kappa=64,128$, and the third ($\theta=1/16$) for $\kappa=256$. The theorem
guarantees success only from $\theta=1/32$ or $1/64$; empirically trials
succeeded when $\theta^2\kappa\le2.2$ and failed when $\theta^2\kappa\ge4$.
On paths with $\kappa\approx2$ and $n=4,\dots,128$, filtered uniform grids
used exactly $4(\lfloor\sqrt{n/4}\rfloor+1)+1$ labels ($9$ to $25$),
matching the separable form of Proposition~\ref{exp:prop:sharp}; graded grids
used $9$ to $19$ ($\theta=1/8$) and $9$ to $15$ ($\theta=1/4$).

\paragraph{Exact output (E4).} On twenty instances with $n\le6$ the exact
wrapper certified $18$ optima, all equal to an independent enumeration of
integer labels and active faces, including three instances with two isolated
optima. The two instances with a segment of optima stopped at the time limit
with certified gaps near $2^{-11}$, while the height rule needs about
$2^{-53}$. The number of stages tracks the required number of gap bits
$\log_2(VW)$: $67$ to $144$ stages for $49$ to $120$ bits, and $275$ and $534$
stages for planted instances needing $208$ and $342$ bits.

\paragraph{Localized acceptance (S1).} On $30$ random mixed-integer instances
($n=4,\dots,16$), the test of Proposition~\ref{exp:prop:local}, applied to the
replayed history of an ordinary run, certified the optimum of $29$ within at
most four stages and $0.031$\,s, whereas the height rule used between $1$ and
$542$ stages. Twelve of the $29$ were certified by the minorant on the full
box; seventeen needed the localization. All $29$ values agreed with the
height rule, and the twelve with $n\le6$ with the enumeration. The remaining
instance has two optimal vertices.

\paragraph{A numerical solver (E5).} Table~\ref{exp:tab:e5} compares SCIP 10.0
(one thread, epigraph formulation, absolute gap $10^{-6}$, $20$\,s limit)
with the default certified solver (absolute gap $10^{-6}$) on planted
instances with $\kappa\approx4.4$. SCIP's incumbents were optimal to about
$10^{-14}$ when evaluated exactly and its dual bounds were valid, but all
$21$ of its reported primal bounds lay below the true optimum, by $1.8\cdot
10^{-6}$ to $1.6\cdot10^{-5}$, because the epigraph constraint is satisfied
only within tolerance. On paths with $n\ge32$ SCIP did not close the gap in
$20$\,s. This is one benign family, one formulation and default settings,
not a performance comparison.

\begin{table}[ht]\centering\small
\begin{tabular}{lrrrrr}
\toprule
& \multicolumn{2}{c}{SCIP (numerical)} & \multicolumn{3}{c}{certified grid}\\
\cmidrule(lr){2-3}\cmidrule(lr){4-6}
Instances (3 each) & time (s) & dual bound at stop & solve (s) & replay (s) & gap\\
\midrule
path, $n=16$ & $1.06$ & $-2.9\cdot10^{-6}$ & $0.33$ & $0.35$ & $\le10^{-6}$\\
tree, $n=16$ & $0.59$ & $-5.3\cdot10^{-6}$ & $0.26$ & $0.36$ & $\le10^{-6}$\\
band, $n=12$, $p=3$ & $1.20$ & $-4.0\cdot10^{-6}$ & $1.25$ & $1.98$ & $\le10^{-6}$\\
band, $n=8$, $p=4$ & $0.17$ & $-3.8\cdot10^{-6}$ & $5.31$ & $10.02$ & $\le10^{-6}$\\
path, $n=32$ & $20^\dagger$ & $-1.2\cdot10^{-4}$ & $0.97$ & $1.06$ & $\le10^{-6}$\\
path, $n=64$ & $20^\dagger$ & $-2.8\cdot10^{-4}$ & $2.46$ & $3.27$ & $\le10^{-6}$\\
path, $n=128$ & $20^\dagger$ & $-7.5\cdot10^{-4}$ & $6.82$ & $9.45$ & $\le10^{-6}$\\
\bottomrule
\end{tabular}
\caption{Medians over three planted instances ($\OPT=0$). $^\dagger$time
limit reached with the gap open. Certified gaps are exact rationals replayed
by the checker.}
\label{exp:tab:e5}
\end{table}
